\chapter{格式规范}

正文部分的篇幅（包括绪论、结论、图、表和公式等），按本规范排版，
硕士学位论文一般为40-60页，硕士学位论文，一般不少于30篇，
其中，期刊文献不少于20篇，国外文献不少于10篇（外文文献不少于3篇），
以近5年的文献为主。博士学位论文一般为80-120页，
博士学位论文，一般不少于100篇，其中，期刊文献不少于60篇，
国外文献不少于40篇（外文文献不少于30篇），以近5年的文献为主。

\section{页面格式设置}

\begin{figure}[htbp]
    \centering
    % 第一张子图
    \begin{subfigure}[b]{0.6\textwidth}
        \centering
        \includegraphics[scale=0.4]{page_settings}
        \caption{}
        \label{fig:sub-top}
    \end{subfigure}

    % 这里留一个空行，或者写 \\ 来实现上下换行

    % 第二张子图
    \begin{subfigure}[b]{0.6\textwidth}
        \centering
        \includegraphics[scale=0.4]{page_paper}
        \caption{}
        \label{fig:sub-bottom}
    \end{subfigure}
    
    \caption{页面设置标准图}
    \label{fig:page-setting}
\end{figure}

\begin{enumerate}
  \item 纸张和印刷：纸型为 A4（21.0cm  29.7cm）80g 标准，双面印刷。
  \item 页眉：从摘要页到最后页，每页均须有页眉，五号字，居中编排。
  奇数页页眉为相应内容的名称、正文中相应各章的名称，
  偶数页页眉为“西南科技大学博士学位论文”或“西南科技大学硕士学位论文”。
  页眉的文字内容之下划两条横线，线粗 0.5 磅，线长与页面齐宽。
  \item 页脚：建议采用文本编辑软件的“页脚自动生成功能”生成页脚，
  一般编排在相应正文同一页最下部并与正文部分用细线隔开，
  线长约为 1/4 页宽。五号字。
\end{enumerate}

\subsection{正文内容、页码格式}
\begin{enumerate}
  \item 正文：除以上 3 级标题、图题、表题之外，正文文字均采用小四号字。正文第一章至结束中间不空页。
  \item 字距和行距：除特殊说明外，全文一律采用无网格、1.2倍行距，段前段后不空行。
  \item 页码：论文页码的第 1 页从正文开始用阿拉伯数字标注，直至全文结束。正文前的内容（除封面）用罗马数字单独标注页码。页码位于页面底端，对齐方式为“居中”，页码格式为最简单的数字，不带任何其它的符号或信息。
\end{enumerate}

\section{标题格式总结}

标题按一级、二级和三级等分级，对应的编号为：1、1.1 和 1.1.1。
其它标题和编号的编排原则为：下级标题的显目程度不超过上一级，不重复或混淆。
可采用 1）、（1）、a、a）、（a）等格式。

章标题按一级标题编排，节标题按二级标题编排，小节标题按三级标题编排。
编号与标题之间空一格。正文第一章从奇数页起排。

一级标题：宋体加粗，居中，三号字，段前2行，段后1行。

二级标题：1.1 ********，宋体，左对齐顶格，小三号字，段前1行，段后0.5行。

三级标题：1.1.1********，宋体，首行缩进，四号，段前0.5行，段后0行。

图题和表题：采用中文，居中，五号字，段前空 0.5 行，段后0.5行。
